\documentclass[11pt,a4paper]{article}
\usepackage[margin=24mm]{geometry}
\usepackage{amsmath,amssymb,amsthm}
\setlength{\emergencystretch}{2em}
\newtheorem{lemma}{Lemma}
\newtheorem{theorem}[lemma]{Theorem}
\newcommand{\gen}[1]{\langle#1\rangle}
\title{Selected ambient cuts and three prefix orbits}
\author{Human Compression project: L044 continuation}
\date{}
\begin{document}\maketitle
Let $G=\gen{\rho_0,\ldots,\rho_{r-1}}$ act transitively on
$\Omega$, with involutory marks and string commutations. Eligibility
always means crossing two orbits of the actual full ambient omission
$G_j$. The first lemma extends the acquired ambient interval proof
to selected colours inside a larger interval. The subgroup whose
orbits are counted retains \emph{all} colours of that larger interval.

\begin{lemma}[Selected ambient confinement]\label{confine}
Let $I=\{u,\ldots,v\}$, $T=G_I$ and $F\subseteq I$.
Suppose $V$ is a proper $T$-invariant subset of $\Omega$.
For every $j\in F$, suppose at least two ambient-eligible
$j$-edges exist and every such edge has both endpoints in $V$.
Then, including singleton $T$-orbits in the count,
\[
                      2|F|\le |V|-o_T(V).                \tag{1}
\]
The set $F$ need not be an interval, and marks in $I\setminus F$
need not be trivial on $V$.
\end{lemma}
\begin{proof}
If $F$ is empty the assertion is immediate. Otherwise choose two
eligible edges per colour in $F$ on the full vertex set $V$.
The complete acquired selected-colour square-normal-form proof
uses the actual ambient omissions just as it uses full local
omissions: every alternative path avoiding a colour lies in its
ambient omission, and conjugation by any other mark preserves its
orbit partition. Thus long-cycle reduction, incident-square merging,
ordinary slides and global middle moves remain available.
All newly chosen edges remain in $V$ by the all-edge confinement
hypothesis. A middle move for a square of colours $i-1,i+1$
uses $u<i<v$ and hence $\rho_i\in T$, which preserves $V$.
The selected-colour proof explicitly permits this middle colour to
be unselected: retain the square and transport the other selected
edges by $\rho_i$; there are simply no selected $i$-edges to
account for when $i\notin F$.

Choose a square-only selection with minimum square count $q$.
The paid minimum-distance argument therefore gives at most one
square per component in every such minimum selection. Each
component is a tree or is unicyclic. All selected edges lie inside
actual $T$-orbits. We prove that each nontrivial orbit $O\subseteq V$
has a tree component, by checking the full acquired ambient-orbit
exit argument with the changed colour set.

Suppose a minimum selection has no tree in $O$; its selected edge
count is $E_O=|O|$. A square colour is localized to its square's
component: an eligible edge inside another component would have
an alternative selected path avoiding that colour, while an edge
between components would replace a square edge and decrease $q$.
The complete ordinary-slide proof consequently still preserves
component vertex sets. It uses only selected bridge and square
colours, all in $F$, and is independent of missing other colours.

Every full neighbour of a square vertex stays in its component.
Neighbours of a square colour already lie in the square.
If the neighbour's colour commutes with a square colour, transport
that eligible square edge and use the localization just proved.
The only remaining case is square colours $i-1,i+1$ and neighbour
colour $i$. Here $\rho_i\in T$ preserves each actual orbit,
including $O$, even if $i\notin F$. The global move preserves
$E_O=|O|$. Viewed on the transported selection, it removes the
transported square and inserts the original square. Removal splits
the transported component into four trees. An inserted vertex
cannot lie in another component of $O$, since that would join two
squares, nor outside $O$, which $\rho_i$ preserves. All four trees
must be rejoined, since otherwise a tree would remain in $O$ despite
$E_O=|O|$. Thus the component partition inside $O$ is transported
by $\rho_i$, and the required neighbour lies in the original
component. This checks the acquired partition argument with an
unselected middle mark; tree components in other $T$-orbits cannot
interfere.

Because $O$ is proper and $G$ is transitive, a full edge leaves $O$.
Over all minimum selections with $E_O=|O|$, choose a shortest
connector consisting of a selected path from a square to the inside
endpoint of such an edge, followed by that exit edge. Its colour
$x$ is outside $I$, since every mark in $I$ preserves $O$.
The selected path is nonempty by the neighbour assertion.
If its first colour commutes with a square colour, the acquired
ordinary slide shortens the path and preserves the exit. Its
repeated-colour exception is excluded by square-colour localization:
as in the full ambient-orbit proof, that exception would place the
completing eligible square edge outside the original component.
Every colour used by this slide is selected, so missing colours
introduce no new case.

The remaining first colour is $i$, with square colours $i-1,i+1$.
Here $u<i<v$ and $x\notin I$ give $|i-x|\ge2$.
The global middle move therefore transports the exit to an actual
$x$-edge. It remains an exit from $O$, because $\rho_i\in T$
preserves $O$ and its complement. The transported selected path
starts on the retained square after its first edge has disappeared;
it gives a shorter connector, still with $E_O=|O|$, a contradiction.

Thus every nontrivial $T$-orbit in $V$ has a tree. Each singleton
is an isolated tree as well. Trees contribute one to vertices
minus edges and unicyclic components contribute zero. Summing
gives (1). The proof used confinement of every eligible selected
colour edge, not merely the original choices.
\end{proof}

\begin{theorem}[Three moving orbits]\label{three}
Retain the complete critical singleton and small-support budget
hypotheses. Write $P=G_{<g}$, $Y=\operatorname{supp}(P)$,
$H=G_{\{2,\ldots,g-1\}}$, $C=a^P$ and $b=a^{\rho_0}$,
where $4\le g<r$, $n=2r\ge14$, $P$ is intransitive on
$Y$, $|Y|\le2g+1$, and every colour $2,\ldots,g-1$ has
two ambient-eligible edges. Then $H$ cannot have exactly two
fixed points and three nontrivial orbits in $Y$.
\end{theorem}
\begin{proof}
Suppose it does. The acquired budget gives $|Y|=2g+1$,
fixed points exactly $a,b$, and three nontrivial orbits whose
total size is $2g-1$. The centralizing involution $\rho_0$
acts on the fixed set as $(a\ b)$ and permutes those three
orbits equivariantly.

Suppose first that it exchanges two orbits $U_1,U_2$ of equal
size. The third orbit $W$ has odd size and is preserved, so
$\rho_0$ fixes it pointwise. Strict growth makes $|C|\ge g+1>2$,
so $C$ meets a moving $H$-orbit. If it meets either $U_i$,
it contains both. Intransitivity of $P$ then forces
$C=\{a,b\}\cup U_1\cup U_2$, the support of $\rho_0$.
Since $\rho_1$ preserves $C$ and every later mark centralizes
$\rho_0$, this would make $G$ intransitive.

Otherwise $C=\{a,b\}\cup W$ and its complement in $Y$ is
$E=U_1\cup U_2$. These are the two $P$-orbits. The involution
$\rho_0$ is free on $E$ and fixes $W$. Strict growth gives
$x=b^{\rho_1}\in W$, so the first-prefix anchor orbit is
$\{a,b,x\}$. The actual natural singleton stabilizer and full
IP give $|P_1|=6$, for $P_1=\gen{\rho_0,\rho_1}$.
The relation $(\rho_0\rho_1)^3=1$ makes $P_1$ fix
$W\setminus\{x\}$ pointwise. Its orbits in $E$ have sizes
two or six, and these are all its ambient orbits of those sizes.
All marks of index at least three centralize $P_1$ and preserve
their union $E$. The first three marks preserve $E$ as well,
since $\rho_2\in H$. This again contradicts transitivity.

Thus $\rho_0$ preserves each of the three $H$-orbits.
On each it is either free or the identity. At least one is fixed
pointwise, since their total size is odd. At least one is free,
since otherwise $\rho_0=(a\ b)$ would be odd. Write the
free orbits as $O_1,\ldots,O_k$, where $k=1$ or $2$ and
$|O_i|=2q_i$. Let $V$ be the union of the other $t=3-k$
orbits, of total size $M$. Thus
\[
                  2\sum_i q_i+M=2g-1.                  \tag{2}
\]

Assign to one free orbit every colour in $I=\{2,\ldots,g-1\}$
which has an ambient-eligible edge in any free orbit. Let $F_i$
be the disjoint assigned sets, and let $F$ consist of the remaining
colours. The acquired internal-pair forest argument gives
\[
                         |F_i|\le q_i-1.               \tag{3}
\]
For clarity, its first part needs no terminal-interval hypothesis:
project one chosen eligible edge of each assigned colour to the
$\rho_0$-pairs in $O_i$. Its endpoints are distinct pairs, since
$\rho_0\in G_j$. An alternate projected path avoiding colour
$j$ lifts using the other chosen colours and $\rho_0$, all in
$G_j$, and contradicts eligibility. Hence these distinct-coloured
projected edges form a forest on $q_i$ vertices, proving (3).
All marks in the actual interval $H$ preserve $O_i$ and its pairs;
no projected IP or stronger equality classification is invoked.

For every $j\in F$, all its ambient-eligible edges lie in $V$:
none lies in a free orbit by the definition of $F$, and the other
points are fixed by $H$. The set $V$ is proper and $H$-invariant,
with exactly $t$ actual $H$-orbits. Lemma~\ref{confine}, with
the full interval $I$ and its selected subset $F$, gives
\[
                            2|F|\le M-t.                \tag{4}
\]
This also holds when $F$ is empty. Summing the disjoint colour
assignments and using (2)--(4), we obtain
\[
\begin{split}
g-2&=\sum_i |F_i|+|F|\\
&\le\sum_i(q_i-1)+\frac{M-t}{2}
 =\frac{2g-1}{2}-k-\frac{3-k}{2}
 =g-2-\frac{k}{2}<g-2,
\end{split}
\]
the final contradiction.
\end{proof}

\paragraph{Attribution and exact boundary.}
The selected-colour normal form, full ambient-orbit localization and
interval exit proof are acquired arguments; their historical global
middle move retains the root collaborator's disclosed contribution
and distinct review. Lemma~\ref{confine} combines their exact scopes
by retaining the full interval subgroup while selecting fewer colours.
The confinement and pair-forest budget in Theorem~\ref{three} are
this continuation. The source setting is Cameron--Fernandes--Leemans--Mixer,
\emph{Highest rank of a polytope for $A_n$}, Propositions~5.28--5.32.
This excludes the stated third-orbit configuration. It does not itself
assemble the entire singleton route or settle the separate imprimitive,
all-paired, finite-base and lower-existence obligations of the principal
L044 contract. No finite census, new formal declaration or whole-original
credit is added.
\end{document}
